\typeout{NT FILE chapter-conclusion-retrospective.tex}%
\chapter{Conclusion and Retrospective}
\label{cha:conclusion-retrospective}
This chapter summarizes the work conducted during the internship and its
continuation, highlighting the key achievements, challenges faced, and lessons
learned. It also provides a retrospective analysis of the project, reflecting
on the development process and the overall impact of the work.

\section{Summary of Achievements}

The work carried out during the internship and its continuation resulted in two
complementary software systems, designed and implemented in C++ over the
course of the project.

The \gls{sdr-control-system} provides a higher-level abstraction over the
\gls{DTAPI}, exposing device management, channel configuration, and signal
transmission functionality through a \gls{gRPC} interface. This addresses
the original problem identified at the outset of the internship: that
configuring \gls{SDR} hardware directly through the \gls{DTAPI} required
detailed knowledge of low-level hardware concerns, such as FIFO buffer sizing
and modulation parameter composition, which the developed system abstracts
away behind a clean, strongly-typed, and remotely accessible interface. The
client-server architecture, together with \gls{gRPC}'s reflection
capabilities, additionally enables interactive use through a \gls{CLI}
client, lowering the barrier for engineers unfamiliar with the underlying
hardware to issue commands and inspect available functionality.

The \gls{rs232-interface} provides a lightweight connectivity layer over
\gls{RS232}, using \gls{PPP} to turn the serial connection into a standard
point-to-point IP channel. This allowed the same gRPC interface used by the
\gls{sdr-control-system} to be reused over the serial path, eliminating the
need for a separate bespoke command protocol and keeping the architecture more
uniform across deployment scenarios.

Both systems share a consistent set of design principles established during
the work: the use of errors as values rather than exceptions, the adoption of
CMake as the build system, and a preference for minimal, composable
abstractions over heavier frameworks where the simpler approach was
sufficient. Beyond the two systems themselves, the work also included their
integration with existing proprietary systems at \gls{GMV}, as described in
\autoref{cha:integration-testing}, validated through a combination of unit
testing, code coverage analysis, and the external tester component developed
to simulate measurement equipment in the prototype environment.

\section{Challenges Faced}

Several challenges were encountered throughout the internship, spanning both
technical implementation details and the evolving scope of the project
itself.

A recurring challenge concerned the \gls{DTAPI} documentation itself.
In practice, some API calls exhibited side effects that were not clearly
described, and certain behavioural details relevant to integration were either
incomplete or not explicitly documented. This made debugging and system
integration more time-consuming, as part of the expected behaviour had to be
inferred empirically through controlled tests against real hardware.

Hardware-level challenges were also significant, particularly around serial
port configuration. The WCH CH382 PCIe card initially relied on the generic
\texttt{8250} driver, which proved insufficient for reliable operation and
was replaced with the vendor-supplied \texttt{ch35\_38x\_linux} driver,
exposing the expected \texttt{/dev/ttyWCH0--3} device nodes. Further issues
were traced to incorrect \texttt{termios} configuration, including the
handling of the \texttt{VMIN} parameter. Another quite relevant issue was the
use of two \gls{DTE} devices connected back-to-back without a null modem adapter, which, in simple terms,
meant a device transmitted data into the other's transmit line instead of its receive line. 
This was resolved by using a null modem adapter, which correctly cross-connected
the transmit and receive lines between the two devices.

\section{Conclusion}

The internship and its continuation achieved their core objectives: a
\gls{gRPC}-based abstraction over the \gls{DTAPI} was designed and
implemented, providing a programmatic and remotely accessible interface for
\gls{SDR}-based \gls{GNSS} signal transmission, and a complementary
\gls{RS232}-based connectivity layer was developed to support integration with
external equipment. By carrying \gls{PPP} over the serial link and reusing the
same gRPC contract over the resulting IP channel, the final system addressed
the gap identified in \autoref{cha:state-of-the-art} between low-level,
hardware-specific vendor APIs and the higher-level, programmatic control
required for automated and remotely operated test workflows.

Beyond the specific technical deliverables, the work provided practical
exposure to the realities of developing software against real, professional
hardware. The need to adapt the original work plan in response to a changed
protocol requirement, an evolved transport architecture, and an expanded
integration scope also reflected a broader lesson about software engineering in
a professional context: requirements and boundaries are rarely fixed in
advance, and designing systems with clear interfaces pays off when that
flexibility is inevitably required.

Overall, the work developed during the internship and its continuation
represents a meaningful and reusable contribution within \gls{GMV}, providing a
foundation
that can be extended to support additional \gls{SDR} hardware features,
additional serial deployment scenarios, or further integration with other
systems, beyond what was required within the scope of the original internship.
